\documentclass[border=10pt,tikz]{standalone}
\usepackage{fontspec}
\setmainfont{Apple SD Gothic Neo}
\setsansfont{Apple SD Gothic Neo}
\usepackage{tikz}
\usetikzlibrary{positioning,arrows.meta,shapes,fit,backgrounds,calc,decorations.pathreplacing}
\begin{document}
% _design.tex — shared TikZ design system for all blog diagrams
% Inputs nothing on its own — meant to be \input{} from each diagram .tex
% AFTER \begin{document}, BEFORE \begin{tikzpicture}.
%
% Provides a unified visual language across GoF / DSA / EMC++ / Beautiful series:
%   - Colors: muted, accessible, color-coded by role
%   - Node styles: consistent sizing, rounded, subtle borders
%   - Edge styles: UML-style inheritance, aggregation, plus dashed alternates
%
% Usage:
%   \documentclass[border=10pt,tikz]{standalone}
%   \usepackage{tikz}
%   \usetikzlibrary{positioning,arrows.meta,shapes,fit,backgrounds}
%   \begin{document}
%   \input{../_design.tex}     % adjust relative path
%   \begin{tikzpicture}[blog]  % apply 'blog' style
%   ...

\definecolor{absfill}{RGB}{220,232,247}   % cool blue
\definecolor{absbord}{RGB}{72,118,182}
\definecolor{confill}{RGB}{249,222,222}   % warm red
\definecolor{conbord}{RGB}{197,93,93}
\definecolor{impfill}{RGB}{223,238,222}   % soft green
\definecolor{impbord}{RGB}{99,154,93}
\definecolor{clifill}{RGB}{251,239,205}   % pale yellow
\definecolor{clibord}{RGB}{200,167,72}
\definecolor{hifill} {RGB}{253,224,200}   % accent orange
\definecolor{hibord} {RGB}{210,135,68}
\definecolor{nodefill}{RGB}{250,250,250}
\definecolor{nodebord}{RGB}{120,120,120}
\definecolor{edgec}  {RGB}{80,80,80}
\definecolor{edgesc} {RGB}{145,145,145}
\definecolor{groupbord}{RGB}{200,200,200}

\tikzset{
  blog/.style={
    font=\sffamily\small,
    every node/.style={font=\sffamily\small}
  },
  % --- node roles ---
  cls/.style={
    rectangle, rounded corners=2pt,
    draw=nodebord, line width=0.4pt,
    minimum height=8mm, minimum width=28mm,
    inner sep=3pt, fill=nodefill, align=center
  },
  abs/.style={cls, draw=absbord, fill=absfill, font=\sffamily\itshape\small},
  con/.style={cls, draw=conbord, fill=confill},
  imp/.style={cls, draw=impbord, fill=impfill},
  cli/.style={cls, draw=clibord, fill=clifill},
  hi/.style={cls, draw=hibord, fill=hifill, font=\sffamily\bfseries\small},
  % --- circle (for trees / graph nodes) ---
  vx/.style={
    circle, draw=absbord, line width=0.4pt,
    minimum size=8mm, fill=absfill, inner sep=0pt
  },
  vxc/.style={vx, draw=conbord, fill=confill},
  vxh/.style={vx, draw=hibord, fill=hifill, font=\sffamily\bfseries},
  % --- decision diamond ---
  q/.style={
    diamond, draw=clibord, line width=0.4pt, fill=clifill,
    inner sep=2pt, aspect=2.4, align=center
  },
  % --- edges ---
  inh/.style={                                  % UML inheritance
    -{Triangle[length=3mm,width=3mm,open]},
    draw=edgec, line width=0.5pt
  },
  agg/.style={                                  % UML aggregation
    {Diamond[length=2.5mm,width=2.5mm,open]}-{Stealth[length=2.2mm]},
    draw=edgec, line width=0.45pt
  },
  uses/.style={                                 % directed reference
    -{Stealth[length=2.2mm]}, draw=edgec, line width=0.5pt
  },
  arr/.style={                                  % plain arrow (graphs/trees)
    -{Stealth[length=2mm]}, draw=edgec, line width=0.45pt
  },
  edge/.style={                                 % undirected (trees)
    draw=edgec, line width=0.45pt
  },
  alt/.style={                                  % dashed: similar/alt
    -{Stealth[length=2mm]}, dashed, draw=edgesc, line width=0.45pt
  },
  ret/.style={                                  % dashed creates/returns
    -{Stealth[length=2mm]}, dashed, draw=edgesc, line width=0.45pt
  },
  thr/.style={                                  % threading link (red dashed)
    ->, dashed, draw=conbord, line width=0.5pt
  },
  % --- group / region ---
  group/.style={
    rectangle, rounded corners=4pt,
    draw=groupbord, line width=0.4pt,
    inner sep=10pt, fill=black!2
  },
  glabel/.style={font=\sffamily\bfseries\small, anchor=north west, inner sep=2pt}
}

% Korean font convenience macro — included by every Hangul diagram
\providecommand{\KOR}{}

\begin{tikzpicture}[blog, muxb/.style={cls, minimum width=8mm, text width=7mm, minimum height=12mm, inner sep=1pt, font=\sffamily\scriptsize}, ffb/.style={con, minimum width=14mm, text width=12mm, minimum height=12mm, inner sep=1pt}]
\node[font=\sffamily\bfseries\large] at (6.4,3.6) {Scan chain — 플립플롭을 한 줄로 꿰기};
\node[neut, minimum width=150mm, text width=145mm, minimum height=9mm] (logic) at (6.4,2.2) {조합 논리 (기능 모드에서 각 플립플롭의 입력을 계산)};
\node[muxb] (m0) at (0.0,0) {MUX};
\node[ffb] (f0) at (1.5,0) {\textbf{FF0}\\[2pt]\scriptsize D \quad Q};
\draw[arr] (m0.east) -- (f0.west);
\draw[arr] (logic.south -| m0.north) -- (m0.north);
\draw[alt] (f0.north) -- (f0.north |- logic.south);
\draw (m0.south) -- (0.0,-1.4);
\node[muxb] (m1) at (3.6,0) {MUX};
\node[ffb] (f1) at (5.1,0) {\textbf{FF1}\\[2pt]\scriptsize D \quad Q};
\draw[arr] (m1.east) -- (f1.west);
\draw[arr] (logic.south -| m1.north) -- (m1.north);
\draw[alt] (f1.north) -- (f1.north |- logic.south);
\draw (m1.south) -- (3.6,-1.4);
\node[muxb] (m2) at (7.2,0) {MUX};
\node[ffb] (f2) at (8.7,0) {\textbf{FF2}\\[2pt]\scriptsize D \quad Q};
\draw[arr] (m2.east) -- (f2.west);
\draw[arr] (logic.south -| m2.north) -- (m2.north);
\draw[alt] (f2.north) -- (f2.north |- logic.south);
\draw (m2.south) -- (7.2,-1.4);
\node[muxb] (m3) at (10.8,0) {MUX};
\node[ffb] (f3) at (12.3,0) {\textbf{FF3}\\[2pt]\scriptsize D \quad Q};
\draw[arr] (m3.east) -- (f3.west);
\draw[arr] (logic.south -| m3.north) -- (m3.north);
\draw[alt] (f3.north) -- (f3.north |- logic.south);
\draw (m3.south) -- (10.8,-1.4);
\draw[flow] (f0.east) -- (m1.west);
\draw[flow] (f1.east) -- (m2.west);
\draw[flow] (f2.east) -- (m3.west);
\node[hi, minimum width=16mm] (si) at (-2.0,0) {scan\_in};
\draw[flow] (si.east) -- (m0.west);
\node[hi, minimum width=16mm] (so) at (14.6,0) {scan\_out};
\draw[flow] (f3.east) -- (so.west);
\draw (0,-1.4) -- (10.8,-1.4);
\node[caption, anchor=north] at (5.4,-1.6) {SE (scan enable) — 0이면 기능 경로, 1이면 굵은 화살표의 시프트 경로};
\node[caption, anchor=north] at (5.4,-2.3) {점선: 기능 모드에서 각 플립플롭의 출력이 다시 조합 논리로 들어감};
\end{tikzpicture}
\end{document}
